% Generated draft: requires mathematical and source review.
\documentclass[11pt]{article}
\usepackage[margin=25mm]{geometry}
\usepackage{fontspec}
\setmainfont{Latin Modern Roman}
\usepackage{amsmath,amssymb,mathtools,unicode-math}
\setmathfont{Latin Modern Math}
\usepackage{enumitem,needspace,xurl,hyperref}
\setlength{\parindent}{0pt}
\setlength{\parskip}{5pt}
\setlength{\emergencystretch}{3em}
\newcommand{\N}{\mathbb{N}}
\newcommand{\Z}{\mathbb{Z}}
\newcommand{\Q}{\mathbb{Q}}
\newcommand{\R}{\mathbb{R}}
\newcommand{\C}{\mathbb{C}}
\newcommand{\F}{\mathbb{F}}
\newcommand{\A}{\mathbb{A}}
\newcommand{\PP}{\mathbb{P}}
\newcommand{\E}{\mathbb{E}}
\newcommand{\eps}{\varepsilon}
\newcommand{\dd}{\,\mathrm{d}}
\begin{document}
\section*{UnsolvedMath ID: 701491}
\textbf{Source problem code:} catalogue entry
\subsection{Definitions and mathematical statement}
For a finite group \(H\), define its spectrum by
\[
\omega(H)=\{\operatorname{ord}(h):h\in H\}.
\]
A section of \(H\) is a quotient \(K/N\), where \(K\leq H\) and \(N\trianglelefteq K\). A proper section means such a quotient with \(|K/N|<|H|\). For a finite set \(\Omega\subseteq\Z_{>0}\), call a finite group \(H\) \(\Omega\)-critical if
\[
\omega(H)=\Omega
\quad\text{and}\quad
\forall K\leq H\;\forall N\trianglelefteq K,
\qquad |K/N|<|H|\Longrightarrow\omega(K/N)\ne\Omega.
\]
Let
\[
\mathcal C(\Omega)=\{[H]:H\text{ is a finite }\Omega\text{-critical group}\},
\]
where \([H]\) denotes the isomorphism class of \(H\). Counting and classification are interpreted up to isomorphism; the supplied statement does not explicitly specify this convention. The finite-group domain follows the source's definition of spectrum. A finite simple group is understood to be nontrivial, with no restriction to nonabelian groups.

\begin{enumerate}
\item[(a)] Does there exist an integer \(n\geq1\) such that, for every finite simple group \(G\),
\[
|\mathcal C(\omega(G))|<n?
\]
The same \(n\) must work for every finite simple group \(G\).
\item[(b)] For every finite simple group \(G\), determine \(\mathcal C(\omega(G))\); that is, classify up to isomorphism all finite \(\omega(G)\)-critical groups.
\end{enumerate}
These are the two questions attributed to V. D. Mazurov. The supplied registry labels the problem open; that label is supplied source metadata and is not an independently verified mathematical or literature assessment.
\subsection{Short English statement}
A finite group is spectrum-critical if it has the prescribed element orders and every proper subgroup quotient has a different spectrum. Mazurov asks whether the number of such groups, counted up to isomorphism, is uniformly bounded for spectra of finite simple groups, and asks for their classification for each such spectrum.
\subsection{Sources}
\begin{itemize}
\item Supplied local UnsolvedMath registry record: permanent ID \(701491\), exact source code \texttt{catalogue entry}, title \emph{Kourovka Notebook Problem 18.62}. The supplied status is \texttt{open}; the record was updated on \texttt{2026-10-16}.
\item Attributed source: \emph{The Kourovka Notebook (Unsolved Problems in Group Theory)}, original problem 18.62, 18th issue (2014); also identified in the supplied background with the 21st issue edition (2026), \texttt{arXiv:1401.0300v48}. Proposed by V. D. Mazurov in 2014. The supplied collection description names E. I. Khukhro and V. D. Mazurov as editors.
\item The source attribution and edition information are supplied registry data. No independent source consultation or verification is asserted, and no separate literature-check date is supplied.
\end{itemize}
\end{document}
